\section{Gott der Eifersüchtige}
\begin{block}[Einführung]
    Donnerstag Abend, sass ich so am Computer fragte Gott, welchen Namen ich denn für den nächsten Sonntag nehmen soll. Ich ging nochmals die Liste durch und es sprang mir der Name Gott der Eifersüchtige ins Auge. Meine Reaktion: \enquote{Super, das mach ich!} Die Datei bereitgemacht, den ersten Bibelvers \biblezit{IIMos}{20:5} gelesen, wie um alles in der Welt erkläre ich das nun? Ich hatte keine Ahnung, was ich jetzt damit anfangen soll. Der Erste Teil des Verses, war ja noch einfach. Gott will nicht, dass wir neben ihm jedamnd anderen anbeten. Der zweite Teil vom Vers, wird es schwieriger. Wie ich euch schon oft gesagt habe, oder wenigstens versucht habe zu sagen, ist der Gott vom AT der gleiche wie der Gott vom NT und der Gott dem wir heute hier in diesem Gottesdienst danken und ihm die Ehre erweisen.

    Er ist der Gott, der seinen Sohn für unsere Rettung töten liess. Er ist die Liebe. Wie hier auf der Liste steht, Gott der Allmächtige, Gott der Arzt, Herr der Heerscharen usw. Jetzt steht da plötzlich \begin{bibelbox}{SCHL}{IIMos}{20:5}
    \dots der die Schuld der Väter heimsucht an den Kindern bis in das dritte und vierte Glied derer, die mich hassen.
    \end{bibelbox}
    Kann ich jetzt einfach über den Namen \betonung[b, p]{Gott der Eifersüchtige} etwas erzählen und diesen zweiten Teil ignorieren? Lesen wir doch noch drei andere Verse die Gott als Eifersüchiger Gott ansprechen.
    \biblezit{VMos}{4:24} dann \biblezit{VMos}{5:9} und \biblezit{VMos}{6:15}. Puh, das sind harte aussagen. Die einfachste Erklärung ist, das galt nur den Juden und uns Christen gilt das nicht mehr. So, jetzt sind wir alle zufrieden oder?

    Ich glaube das ist zu einfach. Wie schon betont ist Gott, \betonung[b, p]{ich bin}.  Das heisst er war, er ist und er wird immer sein. Jesus bezieht dieses \betonung[b, p]{ich bin} auf sich selber. \biblezit{Joh}{15:1}. 

    Als Jesus sah, dass der Tempel als Markthalle benutzt wurde, trieb er sie alle mit einer Geissel aus dem Tempel. Als die Jünger das mitbekamen, kam ihnen den \biblezit{Ps}{69:10} in den Sinn. \biblezit{Joh}{2:17}. Auch Jesus eiferte um seinen Vater.

    In \biblezit{Jak}{4:5} schreibt Jakobus, dass der Heilige Geist in uns ein eifersüchtiges Verlangen hat uns in der Gemeinde Jesus zu halten. Er hilft uns uns, er unterstützt uns und lässt uns nicht alleine. Er will nicht, dass uns der Widersacher oder Welt uns wieder von Jesus wegführt. Das ist das schlechte Gewissen,  er bewacht uns eifersüchtig.

    Die definition von Eifersucht ist:
    \begin{quote}
        Eifersucht ist das starke Gefühl, etwas Wertvolles, das einem gehört oder zusteht, an jemand anderen zu verlieren oder mit jemand anderem teilen zu müssen.        
    \end{quote}
    Jesus ist für dich und mich am Kreuz gestorben. Kurz bevor er in den Himmel aufgefahren ist, sagt er zu den Jüngern \biblezit{Apg}{1:8}. Zu Pfingsten wurde der Heilige Geist ausgeschüttet. Als wir Jesus Christus bezeugten und in ihm Wiedergeboren wurden, bekamen wir auch den Heilige Geist, von dem wir eifersüchtig überwacht werden, dass uns niemand mehr von Jesus wegreissen kann. Gott ist nicht eifersüchtig auf dich, weil du es so schön hier hast. Gott wacht eifersüchtig auf dich weil du ihm so wertvoll bist.

    Das ist ein Grund, weshalb wir ein Gottgefälliges Leben führen sollen. Nicht weil wir dann besser vor ihm stehen, sondern weil jeder einzelne von uns etwas Wertvolles für ihn sind.

    So auch schon im Alten Testament. Er hat sein Volk gewarnt. Für ihn ist sein Volk Wertvoll. Gott bezeichnet sein Volk als seinen Augapfel den er eifersüchtig hütet. Aber seine Eifersucht ist nicht besitzergreifend oder erdrückend. Wer Gott meidet den lässt er  gehen. Gott ist aber enttäuscht, weil er etwas Wertvolles verloren hat.

\end{block}